\chapter{决策树}
\label{chap:trees}
\chaptergoal{能够把一棵小树翻译成判断规则，完整计算分类和回归的候选划分；能够说明贪心生长为什么不保证整棵树最优；理解深度、叶样本数与代价复杂度剪枝的不同作用；利用实际训练结果解释预测路径、泛化表现和适用范围。}

前面的逻辑回归通过一个得分函数形成类别判断，近邻方法则依据附近训练样本作出预测。本章研究第三种组织方式：先问一个容易回答的问题，再根据答案进入下一个问题，直到得到预测。例如，先检查花瓣长度，再检查花瓣宽度，最后确定预测物种。模型要学习的不再只是一组直线系数，还包括用哪个特征提问、阈值在哪里，以及需要问多少层。

本章讨论的是\textbf{以单个数值特征作二元阈值划分的分类与回归树}，即常见的CART式模型。这是scikit-learn树模型的基本形式，不表示所有决策树都只能二叉或只能采用轴平行划分。分类树和回归树共享递归划分结构，但评价划分的目标不同。本章保留必要公式，并用三组小数据把规则、训练目标和模型行为联系起来。

\section{从规则到树结构}
\subsection{节点、分支与预测路径}
决策树是一种按条件逐步划分输入空间、在最终区域给出预测的模型。最上面的\textbf{根节点}接收输入；继续提出条件的节点称为\textbf{内部节点}；不再分支的节点称为\textbf{叶节点}。连接节点的\textbf{分支}表示当前判断的结果。根节点深度为0，每向下经过一条分支，深度增加1。

设某节点检查 $x_j\leq t$，其中 $x_j$ 是第 $j$ 个特征，$t$ 是阈值。满足条件走左分支，不满足走右分支。到达某个叶节点的全部条件构成一条\textbf{预测路径}。路径中的条件是“并且”的关系，而不是任选一个满足即可。不同叶节点对应互不重叠的区域；在边界上使用哪一侧，由 $\leq$ 与 $>$ 的约定明确决定。

分类时，一个叶节点记录进入它的训练样本的类别构成，再据此给出类别。没有样本权重时，通常预测比例最高的类别；若两类比例相同，需要采用实现规定的打破平局规则。回归时，叶节点给出一个数值。使用平方误差准则时，这个数是叶内训练目标的均值。预测新样本不需要知道其目标，只需要沿树中的输入条件找到叶节点。

\worked{先读懂一条人为提出的规则}
下面10条记录是为解释算法而人工构造的，不是鸢尾花或观测目录。
\begin{center}\small
\begin{tabular}{rrrrrrrrrrr}\toprule
$x$&1.0&1.6&2.0&2.4&2.8&3.0&3.4&3.6&4.2&5.0\\
类别& A&A&A&A&A&B&A&B&B&B\\\bottomrule
\end{tabular}\end{center}

暂时由人提出：“若 $x\leq2.5$，预测A；否则预测B。”左叶有4个A，右叶有2个A和4个B。这棵只有一次判断的树称为\textbf{决策树桩}。它在训练记录上错分两个A。对于新输入 $x=3.1$，模型预测B；这并不意味着3.1附近一定不存在A，只表示这条规则把该输入分配给以B为多数类的区域。

现在要区分两个问题：\textbf{如何使用已经给定的树预测}，与\textbf{如何从数据学习这棵树}。前者只是执行条件；后者必须解释为什么采用2.5而不是其他阈值。下一节会发现，2.5只是便于读树的人为示例，并不是这组数据的最优第一次划分。

\subsection{一棵树怎样表示非线性关系}
一个二维判断 $x_1\leq t_1$ 沿竖直线划分平面；在某一侧继续判断 $x_2\leq t_2$，会形成水平边界。多个判断把平面切成若干轴平行矩形，在更高维则是相应的长方体区域。模型可以因此形成弯折或台阶状的决策边界，不要求所有样本遵守一条全局直线。

不过，灵活不等于能用很少的节点表示任意边界。一条倾斜直线对线性分类器很简单，对只作轴平行划分的小树却可能需要多层台阶近似。这说明“哪个模型简单”取决于真实关系怎样与模型的表示方式对应。特征旋转或构造有意义的组合，可能改变树表达同一关系所需要的规模。

\subsection{尺度、顺序与阈值的单位}
树的候选划分依赖特征的排序，而不像近邻那样直接比较不同列的数值距离。对一列进行严格递增变换，原有训练记录的先后顺序不变，因而在理想精确计算下可以保留相同的候选训练划分。由厘米改为米，不会让另一列自动获得更多划分机会，因此树通常不要求为了距离平衡而标准化。

这里不能扩大成“任何变换都不会改变任何预测”。取整会合并不同数值，非单调变换会改变顺序；有限精度和候选并列也可能改变实现结果。即使严格非线性单调变换保留训练分区，以相邻变换值中点作为阈值，映射回原单位后也未必等于原坐标中点，所以落在训练点之间的新输入可能受影响。报告仍须保留阈值的单位，并说明实际变换。

\section{分类树的划分依据}
\subsection{用不纯度描述一个节点的类别混合}
给当前节点记为 $m$，其中有 $n_m$ 个样本、$K$ 个可能类别。第 $k$ 类比例为 $p_{mk}$，满足 $p_{mk}\geq0$ 且 $\sum_kp_{mk}=1$。\textbf{基尼不纯度}定义为
\begin{equation}
G(m)=1-\sum_{k=1}^{K}p_{mk}^{2}.
\label{eq:tree-gini}
\end{equation}
节点只含一个类别时，某个比例为1，其余为0，故 $G=0$。两类各占一半时，$G=1-1/4-1/4=1/2$。一般的最大值是 $1-1/K$，对应 $K$ 类均匀分布。

这个量有一个精确的概率解释：按照节点的经验类别比例，\textbf{独立、有放回地抽取两次类别}，两次相同的概率为 $\sum_kp_{mk}^{2}$，不同的概率就是 $G$。这里明确独立、有放回，是因为有限样本无放回抽取的概率会有修正。不纯度是描述类别混合的统计量，不是某一个新样本被误判的概率。

若只以叶内多数类预测，节点的训练错分比例为 $1-\max_kp_{mk}$。它与Gini不是同一个量。例如类别比例从 $(0.6,0.4)$ 变为 $(0.7,0.3)$ 时，两者都会降低，但降低幅度不同。树使用较敏感的混合指标搜索划分，不等于直接对整棵树的分类准确率做全局优化。

\subsection{候选划分与加权目标}
给定候选特征 $j$ 和阈值 $t$，记录分成左、右子集 $L,R$。划分后目标为
\[
G_{\mathrm{split}}=\frac{n_L}{n_m}G(L)+\frac{n_R}{n_m}G(R),
\]
相应下降量为
\begin{equation}
\Delta G=G(m)-\frac{n_L}{n_m}G(L)-\frac{n_R}{n_m}G(R).
\label{eq:tree-gain}
\end{equation}
对同一个父节点，最小化划分后不纯度与最大化下降量等价。权重表示随机从父节点抽一条记录时进入各子节点的比例；不能把含1条和含99条记录的两个子节点等权相加。

数值特征按值排序后，只有跨过不同训练值才会改变样本归属。因而可以在相邻不同值之间选择代表阈值，通常使用中点。若相邻值都是3.0，就不存在把这两条相同输入按该特征分开的阈值。实际实现还要检查叶样本数等约束；默认最佳划分是在允许的候选中寻找最优，而不是忽略约束无条件切分。

\worked{枚举候选后，最优阈值不是最初的直觉}
对上一节10条记录，父节点比例为 $(0.6,0.4)$，所以 $G(m)=0.48$。比较以下三种情况。

阈值2.6位于2.4和2.8的中点，与人为阈值2.5产生相同训练分区。左侧4A纯净，右侧2A、4B，不纯度为 $4/9$。故加权不纯度为 $0.6\times4/9=4/15$，下降为 $0.48-4/15=0.213333$。

阈值2.9位于2.8和3.0之间。左侧5A纯净，右侧1A、4B，不纯度为 $1-(1/5)^2-(4/5)^2=0.32$。加权后为 $0.5\times0.32=0.16$，下降为 $0.32$。虽然只改变一个样本的归属，这一候选明显优于2.6。

阈值3.5位于3.4和3.6之间。左侧6A、1B，不纯度为 $12/49$，右侧3B纯净。加权不纯度为 $0.7\times12/49=6/35$，下降约 $0.308571$。它与2.9都只错分一个训练样本，但Gini更偏向2.9的分区。

\begin{table}[htbp]\centering\small
\caption{人工10点数据全部候选中点的实际计算。最优下降为0.32，对应阈值2.9。}
\begin{tabular}{rrrr}\toprule
阈值&左/右样本数&划分后加权Gini&下降量\\\midrule
1.3&1/9&0.444444&0.035556\\
1.8&2/8&0.400000&0.080000\\
2.2&3/7&0.342857&0.137143\\
2.6&4/6&0.266667&0.213333\\
2.9&5/5&0.160000&0.320000\\
3.2&6/4&0.316667&0.163333\\
3.5&7/3&0.171429&0.308571\\
3.9&8/2&0.300000&0.180000\\
4.6&9/1&0.400000&0.080000\\\bottomrule
\end{tabular}\end{table}

\fig{v3_ch08_candidates.pdf}{枚举全部候选后的Gini下降量。此图来自人工10点记录，不是鸢尾花训练结果。}

以下短程序直接对应“排序—形成中点—分别统计—比较目标”。它不是为了代替成熟树实现，而是验证第一次划分的机制。
\begin{lstlisting}[caption={枚举一个数值特征的候选阈值}]
import numpy as np
x = np.array([1.,1.6,2.,2.4,2.8,3.,3.4,3.6,4.2,5.])
y = np.array([0,0,0,0,0,1,0,1,1,1])
def gini(labels):
    p = np.bincount(labels, minlength=2) / len(labels)
    return 1. - np.sum(p*p)
values = np.unique(x)
for t in (values[:-1] + values[1:])/2:
    left = x <= t
    after = left.mean()*gini(y[left])
    after += (~left).mean()*gini(y[~left])
    print(t, gini(y)-after)
\end{lstlisting}

\subsection{熵与信息增益}
另一种不纯度为\textbf{信息熵}
\[
H(m)=-\sum_kp_{mk}\log_2p_{mk},\qquad 0\log_20=0.
\]
采用以2为底的对数时单位为比特。两类纯净节点熵为0，均匀节点熵为1。用熵替换式\eqref{eq:tree-gain}中的Gini，所得下降称为\textbf{信息增益}。

本例父节点熵约0.970951。阈值2.6的加权熵为 $0.6H(1/3,2/3)\approx0.550978$，增益约0.419973；阈值2.9的加权熵为 $0.5H(1/5,4/5)\approx0.360964$，增益约0.609987。两者在这组候选中偏向同一划分，但这不是一般定理。不同准则对比例变化的敏感程度不同，可能选出不同树。不能用一次相同排序，推出“选择指标几乎不会影响结果”。

还可以给样本或类别设置权重，此时比例、节点权重和增益都应使用加权量。这样做相当于改变训练时不同错误的代价，不能只修改最终报告中的类别名称。无权重的小例子是理解机制的起点，有权重的实现需要另行核对。

\section{回归树与区域预测}
\subsection{为什么平方误差下的叶值是均值}
当目标是连续数值，一个叶内所有记录共享常数 $c$。平方误差目标为
\[
S(c)=\sum_{i\in m}(y_i-c)^2.
\]
求导得 $S'(c)=2n_mc-2\sum_{i\in m}y_i$，令其为0得到 $c=\bar y_m$；二阶导数为 $2n_m>0$，所以这是唯一最小值。等价地可以把平方和分解为
\[
\sum_{i\in m}(y_i-c)^2=\sum_{i\in m}(y_i-\bar y_m)^2+n_m(c-\bar y_m)^2.
\]
第一项与 $c$ 无关，第二项最小为0，仍得到均值。这就是叶值与第3章损失函数的联系，不是软件任意规定一个平均数。

如果叶内改用绝对误差，最优常数是中位数，可能不唯一。因此“回归树输出叶均值”的说明应限定在本章使用的平方误差树；其他损失会改变叶值和划分目标。目标有测量单位，叶值使用相同单位，平方误差则使用该单位的平方。

\subsection{一次回归划分的完整计算}
定义节点的平均平方偏差
\[
V(m)=\frac1{n_m}\sum_{i\in m}(y_i-\bar y_m)^2.
\]
它使用分母 $n_m$，不是用于无偏估计总体方差时常见的 $n_m-1$。此处的目的，是计算当前训练目标；采用一致分母，才能与整个训练集的平方误差对应。

\worked{比较五个候选阈值}
取人工记录 $x=(1,2,3,4,5,6)$，$y=(1,2,3,7,8,9)$。父节点均值为5，总平方误差为58，故 $V(m)=58/6=29/3$。阈值3.5把前后各三条记录分开，左叶均值2，右叶均值8；每叶平方误差均为2，总和为4，加权平均误差为 $4/6=2/3$。下降量是 $29/3-2/3=9$。

阈值1.5则把目标1单独分出，其余目标均值5.8；平方误差和为38.8，下降为 $(58-38.8)/6=3.2$。其他候选也可逐个计算：
\begin{center}\small
\begin{tabular}{rrrrr}\toprule
阈值&左叶预测&右叶预测&总平方误差&平均误差下降\\\midrule
1.5&1.00&5.80&38.80&3.200\\
2.5&1.50&6.75&21.25&6.125\\
3.5&2.00&8.00&4.00&9.000\\
4.5&3.25&8.50&21.25&6.125\\
5.5&4.20&9.00&38.80&3.200\\\bottomrule
\end{tabular}\end{center}

注意划分条件始终作用于 $x$，不是“看到 $y$ 大于5就放右边”。训练时用 $y$ 评价划分好坏，但预测新样本时必须只使用可获得的输入。若把目标本身用作判断条件，训练误差会很好看，却没有预测能力。

一般地，令 $\mathrm{SSE}(m)$ 表示节点平方误差和，有
\[
n_m\Delta V=\mathrm{SSE}(m)-\mathrm{SSE}(L)-\mathrm{SSE}(R).
\]
因此，对同一父节点最大化加权方差下降，就是最小化分裂后的叶内平方误差总和。这个等价关系使分类树和回归树拥有相似的搜索程序，但它们的具体目标不能混用。

\subsection{分段常数、插值与外推}
树形成区域 $R_1,\ldots,R_M$，预测写成 $\hat y(\mathbf x)=c_m$，其中 $\mathbf x\in R_m$。在同一叶内改变输入，预测不变；跨越阈值后，预测可以突然跳跃。这称为\textbf{分段常数}表示。它适合表达某些分区差异，但并不自动具有物理关系所要求的连续性或单调性。

\fig{v3_ch08_regression.pdf}{人工六点回归树的分段预测与训练范围外行为。图中的两个叶值为2和8；曲线不代表真实住房或物理观测结果。}

对本例，$x=0$仍进入最左叶，预测2；$x=8$进入最右叶，预测8。树不会沿着最后两个训练点的趋势继续增长。多维情形也一样：输入最终仍被路由到已有叶，只是“边界叶”可能取决于其他特征。与线性模型比较时，应分别检查训练范围内的预测和范围外行为，不能因一种模型会给出一个数，就把该数视为可靠外推。

\section{树的生长和停止条件}
\subsection{递归训练的完整过程}
从根节点开始，训练程序检查允许使用的特征和候选阈值，计算划分目标，选择当前最好的合法划分；把记录送入两个子节点；对每个子节点重复这一过程。\textbf{递归}指在子问题中重复同一种操作。若节点已经纯净、没有可用划分，或触发复杂度约束，就不再生长，而是保存叶值。

\textbf{贪心}指每一步依据当前节点的即时改善作选择，不为每个候选预先穷举所有未来子树。树训练依然是在优化一个目标，只是采用离散搜索和局部选择，而不是第3章的普通梯度下降。不能把“没有学习率”解释为“没有训练目标”。阈值、结构和叶值由训练数据确定，最大深度等配置则由建模者通过开发数据选择。

\worked{即时改善最大的第一步，未必得到最好的限深树}
以下人工二维网格含9条记录：
\begin{center}
\begin{tabular}{lccc}\toprule
&$x_2=1$&$x_2=2$&$x_2=3$\\\midrule
$x_1=1$&A&A&A\\
$x_1=2$&A&A&B\\
$x_1=3$&B&B&B\\\bottomrule
\end{tabular}\end{center}
父节点有5A、4B，Gini为 $40/81$。四个根节点候选的下降量为
\[
\begin{array}{c|cccc}
\text{条件}&x_1\leq1.5&x_1\leq2.5&x_2\leq1.5&x_2\leq2.5\\\hline
\Delta G&16/81&25/81&1/81&4/81
\end{array}
\]
故贪心算法首先选 $x_1\leq2.5$。右侧3个B已经纯净；左侧5A、1B，继续用 $x_2\leq2.5$ 分裂，可以得到4A的纯叶，以及1A、1B的混合叶。若最大深度为2，混合叶不能继续细分，至少错分1条记录。

换一种第一步：先用 $x_2\leq2.5$。虽然根节点下降只有 $4/81$，但左侧4A、2B可再用 $x_1\leq2.5$ 完全分开；右侧1A、2B可用 $x_1\leq1.5$ 完全分开。这棵树同样深度不超过2、叶数不超过4，却能全部分对。配套程序分别计算两棵树，得到错分数1与0。

\fig{v3_ch08_greedy_tree.pdf}{九点人工网格的贪心深度2分类树。某个叶仍含两类；另一种根划分可在相同深度预算内获得零训练错分。}

反例说明即时目标与限深整树目标并不等价，不说明贪心在所有数据上都很差。搜索最优整树往往需要比较大量结构，实际算法用计算效率换取局部近似。后面的集成方法能缓解某些单树不稳定性，但不能一般性地保证找回全局最优树，也不应把随机森林描述成贪心算法的数学纠错器。

\subsection{四个停止参数各限制什么}
\begin{center}\small
\begin{tabularx}{\linewidth}{p{0.31\linewidth}X}\toprule
参数&作用及容易混淆的地方\\\midrule
\texttt{max\_depth}&限制根到叶的最长路径；深度2的二叉树至多4个叶，但未必长满4个。\\
\texttt{min\_samples\_split}&当前节点至少有多少记录，才允许尝试继续分裂；不保证每个孩子都有同样多记录。\\
\texttt{min\_samples\_leaf}&候选分裂后每个孩子都必须达到最低记录数；它直接排除过小叶。\\
\texttt{min\_impurity\_decrease}&要求整体加权目标的改善至少达到给定值，不是简单用每个局部节点的未缩放下降作比较。\\\bottomrule
\end{tabularx}\end{center}

最后一项需注意节点权重：在没有样本权重时，sklearn所比较的整体下降是 $(n_m/n)\Delta G$，$n$ 为整棵树的训练样本数。深处只有少量记录的节点，即使局部比例改善明显，对总体目标的贡献也可能很小。九点例子左子节点局部下降为 $1/9$，乘 $6/9$ 后整体贡献为 $2/27$。理解这一点，才能把手算公式与软件参数对应。

默认树通常允许生长得较复杂，但不保证每个叶只有一条样本。重复输入却具有不同目标时，任何单特征阈值都无法把完全相同的输入分开；某些节点也可能没有合法改善。因此“完全生长”是当前搜索与停止规则下不能继续增长，不是每个任务都能获得零误差。

\section{预剪枝与剪枝的基本思想}
\subsection{在生长时限制，与生长后收回}
\textbf{预剪枝}是在生长过程中利用深度、叶样本数或收益阈值，禁止部分分裂。它节省计算并使模型较简单，但可能提前拒绝一个当前改善不大、后续潜力较大的划分。九点反例给出了理解这种风险的背景，但不能据此机械地建议总是保留所有分支，仍需通过验证选择复杂度。

\textbf{后剪枝}先建立一棵较大的树，再把某些内部子树替换为叶节点。它可以在观察一个已长成子树的整体贡献后作决定，但只在原树的子树集合中选择，不能凭剪枝改变最初错误的根划分。后剪枝缓解预先停止的短视，不等于对全部可能树结构进行重新优化。

\subsection{代价复杂度目标及其单位}
对树 $T$，记 $|T|$ 为叶节点数。代价复杂度目标为
\begin{equation}
R_\alpha(T)=R(T)+\alpha|T|,\qquad \alpha\geq0.
\label{eq:tree-ccp}
\end{equation}
在本书与sklearn对应的无权重口径下，$R(T)=\sum_{m\in\mathrm{leaves}}(n_m/n)I(m)$，其中 $I$ 是所采用的节点不纯度。分类可以是Gini，平方误差回归则是节点平均平方偏差。若采用样本权重，应使用相应权重和，而不是原始计数。\texttt{ccp\_alpha}对应这一归一化目标；把 $R$ 改成未除以 $n$ 的误差总和时，数值相同的 $\alpha$ 不再有相同意义。[34]

\worked{给每一个叶增加复杂度代价}
仍用六点回归数据，只比较三棵明确指定的候选树：单叶树 $R=29/3$；阈值3.5的两叶树 $R=2/3$；每个输入独占一叶的六叶树 $R=0$。三者目标分别为
\[
29/3+\alpha,\qquad 2/3+2\alpha,\qquad6\alpha.
\]
当 $\alpha=0.1$，目标约为9.7667、0.8667和0.6，六叶树最低；当 $\alpha=0.2$，目标约9.8667、1.0667和1.2，两叶树最低；当 $\alpha=10$，单叶树为19.6667，两叶树20.6667，六叶树60，单叶树最低。

六叶与两叶候选在 $\alpha=1/6$ 时相等，两叶与单叶在 $\alpha=9$ 时相等。这里只比较三个指定候选，用来理解权衡；真实剪枝路径可能还包含三叶、四叶等子树，不能把这三个候选当作完整软件路径。也不能把训练目标最小的 $\alpha$ 当成最终应选配置：$\alpha$用于产生候选模型，选多少复杂度仍看验证任务。

进一步地，某个子树 $T_m$ 若被替换成一个叶，叶数减少 $|T_m|-1$，训练不纯度增加 $R(m)-R(T_m)$。两种选择相等时的有效代价为
\[
\alpha_{\mathrm{eff}}=\frac{R(m)-R(T_m)}{|T_m|-1}.
\]
它衡量每减少一个叶所付出的拟合代价。软件按相关候选关系构造剪枝路径；读者不必手写完整实现，但应能解释为什么低贡献的细碎分支往往先被收回。

\subsection{实际剪枝结果可以是平台，而非完美U形}
本章鸢尾花实验采用90条训练、30条验证、30条最终保留记录。以完全生长树为起点，预先列出若干 $\alpha$：0、0.005、0.01、0.02、0.04、0.08、0.16。前3个设置得到5叶，0.02和0.04得到4叶，最后两个得到3叶。前5种的验证准确率均为 $29/30\approx0.9667$，最后两种为 $28/30\approx0.9333$。

\fig{v3_ch08_pruning.pdf}{固定开发划分中的剪枝对照。评价基于30条验证记录，因此一条记录就对应约0.0333的准确率变化；平台不应被解释为所有配置完全等价。}

这组结果说明删去一个分支可以保持本次验证准确率，却不能证明它在所有未来数据上都无价值。若在开发阶段比较多个设置，可依据预先规定的简洁性偏好在近似表现下选择较小树；还应结合交叉验证稳定性，而不是读取一个平台上的任意最高小数。

\section{鸢尾花分类与区域住房价值回归}
\subsection{从真实花的记录读出预测路径}
鸢尾花数据的对象、四项厘米测量及三类物种已在第1章介绍。本节仍使用相同字段，不重新赋予数字类别大小含义。分层保留20\%记录作为最终测试后，在120条开发记录内按第4章约定形成90条训练和30条验证；本章不使用最终测试挑选深度。

\begin{lstlisting}[caption={与本章结果一致的浅树训练与路径读取}]
from sklearn.datasets import load_iris
from sklearn.model_selection import train_test_split
from sklearn.tree import DecisionTreeClassifier, export_text
iris = load_iris(as_frame=True)
X_dev, X_hold, y_dev, y_hold = train_test_split(
    iris.data, iris.target, test_size=.2,
    stratify=iris.target, random_state=42)
X_tr, X_val, y_tr, y_val = train_test_split(
    X_dev, y_dev, test_size=.25,
    stratify=y_dev, random_state=7)
tree = DecisionTreeClassifier(max_depth=2, random_state=42)
tree.fit(X_tr, y_tr)
print(export_text(tree, feature_names=list(X_tr.columns)))
print(tree.predict_proba(X_val.iloc[[0]]))
print(tree.decision_path(X_val.iloc[[0]]).indices)
\end{lstlisting}

本次根节点为花瓣长度不大于2.45厘米，满足时预测setosa；否则继续检查花瓣长度是否不大于4.95厘米，分别预测versicolor或virginica。验证样本中的原始行号0具有花瓣长1.4厘米，因此第一步就进入setosa叶。其叶内经验类别比例为 $(1,0,0)$，并不意味着对所有未来这样的花都具有100\%可靠性；它只表示当前训练叶没有其他类别。

\fig{v3_ch08_iris_tree.pdf}{在90条训练记录上拟合的深度2鸢尾花树。图、文本路径与结果表来自同一模型，节点中的样本构成不混用全数据重训结果。}

\begin{center}\small
\begin{tabular}{rrrr}\toprule
最大深度&实际叶数&训练准确率&验证准确率\\\midrule
1&2&0.6667&0.6667\\
2&3&0.9667&0.9333\\
3&4&0.9889&0.9667\\
4&5&1.0000&0.9667\\
6&5&1.0000&0.9667\\
不限制&5&1.0000&0.9667\\\bottomrule
\end{tabular}\end{center}

在这组设置中，验证准确率没有随树进一步长大而下降，而是保持不变。因此不能把“加深必然导致验证下降”写成实验结论。过拟合是一种需要检查的风险，不是每张曲线必须呈现的形状。训练完全正确也不等于已经有足够证据支持跨地区、跨测量设备的识别。

\subsection{回归对照必须明确图中的输入维数}
回归机制的离线必做实验采用本章人工六点数据，完整展示训练、分段预测和外推。第5章还提供300条明确标记为人工生成的响应数据，可以比较同一训练与验证划分下的线性、单树和集成模型。它们用于机制和软件检验，不替代真实房价或物理结论。

区域住房案例仍作为本书的真实公开回归练习。数据需要事先下载或提供CSV，完整获取与运行方法见附录D。使用全部八项输入训练回归树时，不能把“横轴收入、纵轴预测”直接画成唯一阶梯线，因为其他七项输入也会改变路由。课堂可选择两种明确设计：仅用收入中位数另训一个单特征树；或固定其他输入在某组具体值，只沿收入方向画条件切片。前者是较简单的新模型，后者是多特征模型的一个切片，两者都必须在图注写清。

多特征性能比较时，应在完全相同的开发样本、输入集合和指标上与线性回归比较。记录MAE、RMSE、叶数及验证误差分布；结果由本地实际运行写入，不预先声称树优于直线。当前离线发布包没有填入未重新获得的住房数据实验数值，代码会明确报告数据缺失，而不会暗中换成其他回归数据。

\section{过拟合和专业规则解释}
\subsection{训练记录的小变化为什么会改变整棵树}
根节点若有两个近似同好的候选，小幅改变记录就可能改变第一名。一旦根划分不同，后续子节点看到的训练集合也不同，许多后续条件会一起变化。这种结构不稳定性并不要求每个单独阈值都大幅变化；一个较早的差异就可能重新组织整棵树。

因此，解释单棵树时至少区分预测稳定性和结构稳定性。两个结构不同的树可能对大多数新样本给出同样类别；一棵看起来简单的树也可能对阈值附近的样本非常敏感。可以在开发数据内进行重复划分或重采样，统计某些特征被使用的频率、主要预测是否改变，避免只挑选最漂亮的一张树图。

树深 $d$ 对应二叉树至多 $2^d$ 片叶，是表达空间的一个上界，不是实际叶数，也不意味着准确率随深度指数提高。每叶的样本支持、测量噪声、类别重叠和重复输入都会限制实际收益。复杂度选择需要兼顾误差、稳定性、解释成本和任务对错误的容忍，而不仅是节点数量。

\subsection{从路径解释到专业任务}
在天文测光分类中，某条路径可以写成几个星等或颜色条件；这些条件是当前样本、特征和标签下的预测规则，不自动是恒星与星系之间的普适物理边界。测光量与光谱参考类别的获得过程可能不同，输入只能使用部署时可获得的量。相同天体的多次观测要按对象分组，否则树可能识别重复对象而非泛化到新对象。

在风洞声学回归中，一个叶代表某些频率、流速、几何条件组成的区域，叶值是该区域训练记录的平均声压级。应检查同一工况下多频率记录是否被分到两侧、叶样本是否充分，以及条件范围外是否仍被重复赋予边界均值。树给出的分区不应未经物理检验就写成声学机制。

专业报告可以按照“对象—路径—支持—误差—限制”的顺序组织。例如先说明某验证对象的测量值，再逐节点写出判断和单位，报告叶中的训练类别构成，最后检查它实际是否被分错、对测量扰动是否稳定。这样的解释直接回答模型怎样作出本次判断，同时保留证据能支持的范围。第9章会通过组合许多树讨论预测波动，第16章再系统研究特征作用与科学结论。

\chapterpractice{完成10点分类候选的全枚举，核对最佳中点2.9；计算六点回归全部候选；用程序验证九点贪心反例。随后对同一鸢尾花开发划分生成树图、读取一条真实预测路径，比较深度和剪枝。住房练习须明确数据已准备、输入维数与图的含义，未获取的数据不填假结果。}
\chapterexercise
\begin{enumerate}
\item 对一个含5A、3B的节点计算Gini与熵。说明它们与多数类错分率为何不是同一量。
\item 对本章10点记录，分别计算2.6、2.9、3.5三个阈值的类别构成、Gini下降和训练错分数。为什么2.9与3.5错分数相同而Gini不同？
\item 证明平方误差下叶值取均值最优，并用六点例子核对 $n_m\Delta V$ 与误差平方和下降的关系。
\item 写出九点反例中两棵树的完整路径，逐条核对所有样本。为什么后剪枝不能把贪心树的根条件自动改成另一棵树的根条件？
\item 对一个含20条记录的节点，假设整树训练样本为100、局部Gini下降0.1。整体加权下降是多少？与 \texttt{min\_impurity\_decrease} 比较时应使用哪个量？
\item 在三棵指定回归候选中，计算 $\alpha=0.15,0.3,12$ 时的代价复杂度。说明这不是对全部可能子树的完整搜索。
\item 重现鸢尾花深度与剪枝表。验证曲线没有下降时，怎样准确表述结果，而不强行写出“过拟合已发生”？
\item 对包含八个输入的住房回归树，设计一个合法的一维显示方案。说明单特征重训、条件切片和多特征散点图分别表达什么。
\item 选一个天文或物理对象，写出一条带单位的预测路径及其证据限制。分别讨论测量误差、重复记录和外推对这条路径的影响。
\end{enumerate}

\section*{深化练习与方法检验}
本组题强调解释、计算与验证；D编号在本章内独立计数。计算题应写出所用定义和中间步骤，实验题保留实际输出，包括不支持预期的结果。
\begin{enumerate}[label=D\arabic*.]
\item 在本章分类十点数据中，若只检查阈值2.6与3.2，会选择哪一个？为什么“在已检查候选中最好”与“所有允许候选中最好”必须区别？
\item 把所有输入做正线性变换 $x\mapsto1000x$，树的训练候选分组怎样变化？若做一般严格递增非线性变换并仍取变换后相邻值的中点，新查询点的归类是否必然完全不变？
\item 设一个节点四条记录的类别为A、A、B、B，但样本权重为1、1、4、4。分别计算无权重与有权重Gini，说明叶内类别概率怎样解释。
\item 某树把每个训练样本都放进纯叶。列出至少两种“训练全对但并不能说明适用”的原因，并设计一个围绕结构稳定性的开发实验。
\end{enumerate}
